% Compile with: pdflatex lecture_27.tex
% Requires: amsmath, amssymb, amsthm, mathtools, xcolor, mdframed,
%           enumitem, hyperref, pagecolor, microtype
\documentclass[11pt]{article}
\usepackage[margin=1in]{geometry}
\usepackage{amsmath, amssymb, amsthm}
\usepackage{mathtools}
\usepackage{xcolor}
\usepackage{mdframed}
\usepackage{enumitem}
\usepackage{hyperref}
\usepackage{pagecolor}
\usepackage{microtype}

% ── Colour palette ──────────────────────────────────────────
\definecolor{cream}{HTML}{FFFDF5}
\definecolor{defblue}{HTML}{1A4A7A}
\definecolor{thmgreen}{HTML}{1A5C3A}
\definecolor{warnAmber}{HTML}{7A4A00}
\definecolor{dangerRed}{HTML}{7A1A1A}
\definecolor{boxbluebg}{HTML}{EBF3FB}
\definecolor{boxgreenbg}{HTML}{EBF7EF}
\definecolor{boxamberbg}{HTML}{FDF5E6}
\definecolor{boxredbg}{HTML}{FBEBEB}
\pagecolor{cream}

% ── Theorem styles ──────────────────────────────────────────
\newtheoremstyle{defstyle}  {8pt}{8pt}{\normalfont}{0pt}{\bfseries\color{defblue}} {.}{0.5em}{}
\newtheoremstyle{thmstyle}  {8pt}{8pt}{\normalfont}{0pt}{\bfseries\color{thmgreen}}{.}{0.5em}{}
\newtheoremstyle{notstyle}  {8pt}{8pt}{\normalfont}{0pt}{\bfseries\color{warnAmber}}{.}{0.5em}{}
\newtheoremstyle{cautionstyle}{8pt}{8pt}{\normalfont}{0pt}{\bfseries\color{dangerRed}}{.}{0.5em}{}

\theoremstyle{defstyle}
\newtheorem{definition}{Definition}[section]

\theoremstyle{thmstyle}
\newtheorem{theorem}{Theorem}[section]
\newtheorem{corollary}[theorem]{Corollary}

\theoremstyle{notstyle}
\newtheorem*{remark}{Remark}
\newtheorem*{example}{Example}

\theoremstyle{cautionstyle}
\newtheorem*{caution}{Caution}

% ── Framed boxes ────────────────────────────────────────────
\newmdenv[
  backgroundcolor=boxbluebg, linecolor=defblue, linewidth=2pt,
  topline=false, rightline=false, bottomline=false,
  innerleftmargin=10pt, innerrightmargin=10pt,
  innertopmargin=6pt, innerbottommargin=6pt,
  skipabove=8pt, skipbelow=8pt
]{defbox}

\newmdenv[
  backgroundcolor=boxgreenbg, linecolor=thmgreen, linewidth=2pt,
  topline=false, rightline=false, bottomline=false,
  innerleftmargin=10pt, innerrightmargin=10pt,
  innertopmargin=6pt, innerbottommargin=6pt,
  skipabove=8pt, skipbelow=8pt
]{thmbox}

\newmdenv[
  backgroundcolor=boxamberbg, linecolor=warnAmber, linewidth=2pt,
  topline=false, rightline=false, bottomline=false,
  innerleftmargin=10pt, innerrightmargin=10pt,
  innertopmargin=6pt, innerbottommargin=6pt,
  skipabove=8pt, skipbelow=8pt
]{notebox}

\newmdenv[
  backgroundcolor=boxredbg, linecolor=dangerRed, linewidth=2pt,
  topline=false, rightline=false, bottomline=false,
  innerleftmargin=10pt, innerrightmargin=10pt,
  innertopmargin=6pt, innerbottommargin=6pt,
  skipabove=8pt, skipbelow=8pt
]{cautionbox}

% ── Macros ──────────────────────────────────────────────────
\newcommand{\Z}{\mathbb{Z}}
\newcommand{\N}{\mathbb{N}}
\newcommand{\R}{\mathbb{R}}
\newcommand{\id}{\mathrm{id}}
\newcommand{\powset}{\mathcal{P}}

\begin{document}

% ── Title ───────────────────────────────────────────────────
\begin{center}
  {\Large\bfseries Lecture 27 --- Finite Sets}\\[4pt]
  {\normalsize CS 173: Discrete Structures \quad$\cdot$\quad Functions \& Cardinality}\\[2pt]
  {\small\color{gray} University of Illinois at Urbana-Champaign}
\end{center}
\vspace{1em}\hrule\vspace{1.5em}

%=================================================================
\section{Finite Sets and Cardinality Notation}
%=================================================================

\begin{defbox}
\begin{definition}[Finite Set]
We say a set $X$ is \textbf{finite} if
\[
  (\exists\, n \in \N)\;\bigl(|X| = |n|\bigr).
\]
Here $|A| = |B|$ means there exists a bijection $f\colon A \to B$.
\end{definition}
\end{defbox}

\begin{defbox}
\begin{definition}[Cardinality of a Finite Set]
If $|X| = |n|$ and $n \in \N$, then
\[
  |X| \coloneqq n = \{0,\, 1,\, 2,\, \dots,\, n-1\}.
\]
\end{definition}
\end{defbox}

%=================================================================
\section{Injections and Left Inverses}
%=================================================================

\subsection{Setup}

Suppose $f\colon A \to B$ is an injection. We wish to show there exists a function
$g\colon B \to A$ such that
\[
  g \circ f = \id_A.
\]
Such a $g$ is called a \textbf{left inverse} for $f$.

\begin{defbox}
\begin{definition}[Identity Function]
The \textbf{identity function} on a set $A$ is the function $\id_A \colon A \to A$
defined by
\[
  \id_A(a) \coloneqq a \quad \text{for every } a \in A.
\]
\end{definition}
\end{defbox}

\begin{defbox}
\begin{definition}[Left Inverse]
Let $f\colon A \to B$. A function $g\colon B \to A$ is a \textbf{left inverse} for $f$
if
\[
  g \circ f = \id_A,
\]
i.e.\ for every $a \in A$, $g\bigl(f(a)\bigr) = a$.
\end{definition}
\end{defbox}
\newpage

\subsection{Composition with the Identity}
\begin{notebox}
\begin{remark}
For any functions $h\colon A \to A$ and $k\colon A \to A$, composing with the identity
gives
\[
  h = \id_A \circ h, \qquad \id_A \circ k = k.
\]
In particular, $\id_A$ is both a left and right identity for composition of functions
$A \to A$.
\end{remark}
\end{notebox}

%=================================================================
\section{Monomorphisms, Epimorphisms, and Isomorphisms}
%=================================================================

The following terminology packages the left/right inverse characterisations into
standard names.

\begin{defbox}
\begin{definition}[Monomorphism (Monic)]
Let $f\colon A \to B$. We say $f$ is a \textbf{monomorphism} (or $f$ is \textbf{monic})
if and only if
\[
  (\exists\, g\colon B \to A)\;\bigl(g \circ f = \id_A\bigr).
\]
Equivalently, $f$ has a left inverse.
\end{definition}
\end{defbox}

\begin{defbox}
\begin{definition}[Epimorphism (Epic)]
Let $f\colon A \to B$. We say $f$ is an \textbf{epimorphism} (or $f$ is \textbf{epic})
if and only if
\[
  (\exists\, g\colon B \to A)\;\bigl(f \circ g = \id_B\bigr).
\]
Equivalently, $f$ has a right inverse.
\end{definition}
\end{defbox}

\begin{defbox}
\begin{definition}[Isomorphism (of Sets)]
A function $f$ is an \textbf{isomorphism} (of sets) if and only if $f$ is both monic
and epic:
\[
  f \text{ is an isomorphism} \iff f \text{ is monic} \;\wedge\; f \text{ is epic}.
\]
\end{definition}
\end{defbox}

%=================================================================
\section{Main Theorems}
%=================================================================

\begin{thmbox}
\begin{theorem}[Injection $\Rightarrow$ Left Inverse \& Injective Retract]
For any sets $A, B$ and any $f\colon A \to B$: if $f$ is injective, then there exists
$g\colon B \to A$ such that
\[
  g \text{ is surjective} \qquad\text{and}\qquad g \circ f = \id_A.
\]
\end{theorem}
\end{thmbox}

\begin{thmbox}
\begin{corollary}
$|A| \leq |B| \Rightarrow |B| \geq |A|$.
\end{corollary}
\end{thmbox}

\begin{thmbox}
\begin{theorem}[Surjection $\Rightarrow$ Right Inverse \& Injective Section]
For any sets $A, B$ and any $f\colon A \to B$: if $f$ is surjective, then there exists
$g\colon B \to A$ such that
\[
  g \text{ is injective} \qquad\text{and}\qquad f \circ g = \id_B.
\]
\end{theorem}
\end{thmbox}

\begin{thmbox}
\begin{corollary}
$|A| \geq |B| \Rightarrow |B| \leq |A|$.
\end{corollary}
\end{thmbox}

%=================================================================
\section{Trichotomy of Cardinality}
%=================================================================

\begin{thmbox}
\begin{theorem}[Axiom 7 --- Trichotomy of Cardinality]
For any sets $A$ and $B$,
\[
  |A| \leq |B| \;\;\lor\;\; |A| \geq |B|.
\]
Equivalently, exactly one of the following holds:
\[
  \bigl(|A| < |B|\bigr) \;\;\lor\;\; \bigl(|A| > |B|\bigr) \;\;\lor\;\; \bigl(|A| = |B|\bigr).
\]
\end{theorem}
\end{thmbox}

\begin{notebox}
\begin{remark}
Trichotomy says that cardinalities are always comparable.  Note that $|A| < |B|$ means
$|A| \leq |B|$ and $\lnot(|A| \geq |B|)$, i.e.\ there is an injection $A \to B$ but
no injection $B \to A$.  Similarly $|A| > |B|$ and $|A| = |B|$ cover the remaining
cases. This axiom is equivalent to the Axiom of Choice and is not provable from ZF
alone.
\end{remark}
\end{notebox}

%=================================================================
\section{Summary}
%=================================================================

\begin{enumerate}[label=\textbf{\arabic*.}, leftmargin=*, itemsep=4pt]
  \item A set $X$ is \textbf{finite} iff $(\exists\, n \in \N)(|X| = |n|)$, where
        $n = \{0, 1, \dots, n-1\}$.
  \item $|A| = |B|$ iff there exists a bijection $A \to B$.
  \item The \textbf{identity function} $\id_A(a) = a$ is left and right neutral for
        composition.
  \item A \textbf{left inverse} of $f\colon A \to B$ is $g\colon B \to A$ with
        $g \circ f = \id_A$.
  \item \textbf{Mono/Epi/Iso}: $f$ is monic $\iff$ $f$ has a left inverse;
        $f$ is epic $\iff$ $f$ has a right inverse;
        $f$ is an isomorphism $\iff$ $f$ is monic and epic.
  \item Injection $\Rightarrow$ $\exists$ surjective left inverse;
        corollary: $|A| \leq |B| \Rightarrow |B| \geq |A|$.
  \item Surjection $\Rightarrow$ $\exists$ injective right inverse;
        corollary: $|A| \geq |B| \Rightarrow |B| \leq |A|$.
  \item \textbf{Trichotomy} (Axiom~7): for any $A, B$,
        $|A| \leq |B| \lor |A| \geq |B|$; equivalently,
        exactly one of $|A| < |B|$, $|A| > |B|$, $|A| = |B|$ holds.
\end{enumerate}

\begin{notebox}
\begin{remark}[Looking Ahead]
With trichotomy in hand, the Cantor--Bernstein theorem ($|A| \leq |B|$ and
$|B| \leq |A|$ implies $|A| = |B|$) will give us a robust framework for comparing the
sizes of infinite sets.
\end{remark}
\end{notebox}

\end{document}